% !TeX root = ../教材题库_全部条目_转写工作稿.tex
% Source: XJ-B1-C06-TH-P211-Y0813
% 第6章 / 6.1 获取数据的途径与统计的一些基本概念 / 教材 p.211 / PDF 第218页
% OCR未校对：公式与文字均可能有误；未完成数学TeX转写；禁止直接发给学生
% 原题及相邻原解截图（核对用）：../题目原图/XJ-B1-C06-TH-P211-Y0813-01.webp
% 本文件只在转写完成后作为题干使用；OCR 条目尚不是可靠题干。
\begin{bankitem}{p.211 思考（含紧邻说明）}{第6章 / 6.1 获取数据的途径与统计的一些基本概念 / 教材 p.211 / PDF 第218页}{XJ-B1-C06-TH-P211-Y0813}{OCR未校对：公式与文字均可能有误；未完成数学TeX转写；禁止直接发给学生}
\begin{ocrdraft}
[教材p.211 / PDF第218页]
\par
？思考人们借助统计学方法研究实际问题，如我国人口的迁徙是否有规律、
\par
经济预测是否科学、产品质量控制是否达标、网购商品是否畅销等等，首先要做的
\par
工作就是针对问题收集相关数据．那么，如何收集数据呢？
\par
所有统计数据都是来自调查或实验．但是，从使用者的角度看，统计数据主要
\par
来自两条途径：间接来源和直接来源．
\par
1.数据的间接来源
\par
如果与研究内容有关的原信息已经存在，我们只是对这些原信息重新加工、整
\par
理，使之成为统计分析可以使用的数据，就称该原信息为间接来源的数据．具体而
\par
言，常见的间接来源数据有国家各级统计部门公布的统计公报、定期出版的各类统
\par
计年鉴，各类专业期刊、报纸、图书所提供的文献资料，从互联网或图书馆查阅到
\par
的相关资料，等等（如图6．1一1)．间接来源数据又称为二手数据．
\par
d国家数据
\par
0？4
\par
中国统计年鉴
\par
图6．1一1
\end{ocrdraft}
\end{bankitem}
